\documentclass[12pt, titlepage]{article}
\usepackage[margin=1in]{geometry}
\usepackage{graphicx}
\usepackage{float}
\usepackage{amsmath}

%\usepackage{hyperref}

\title{\Huge Final Project Report \\ \normalsize{EGME-3170 - Thermal Systems}}
\author{by \\ Benjamin Paxson \\ Daniel Bourguignon \\ Simon Ross}
\date{\footnotesize{Due} \\ \normalsize{April 28, 2026}}


%\setlength{\topmargin}{-.25in}
%\setlength{\oddsidemargin}{.25in}
%\setlength{\evensidemargin}{.25in}
%\setlength{\textwidth}{6.0in}
%\setlength{\textheight}{9.0in}
%\renewcommand{\baselinestretch}{1.4}
%\usepackage{indentfirst} %Maybe keep this maybe not

\begin{document}
\maketitle

%FIXME remove comments

%FIXME improve abstract
\begin{abstract}
Abstract here - we need contributions from each member somewhere


\end{abstract}

\newpage
\tableofcontents
\listoffigures
\listoftables

\newpage
\section{Introduction}

\section{Assumptions}
\section{Calculations}
\section{Results}
\section{Conclusions}
\newpage
\section{Appendices}

\subsection{part1.m}
\(
\textbf{Fin Efficiencies}\\
\\
m = \sqrt{\frac{2h}{kt}}
\text{\quad - Constant for calculating efficiency}\\
\\
L_c = L + \frac{t}{2}
\text{\quad - Corrected fin length (meters)}\\
\\
\eta_{fin} = \frac{\text{tanh}(mL_c)}{mL_c}
\text{\quad - Single fin efficiency}\\
\\
\\
\textbf{Areas}\\
\\
A_b = wt
\text{\quad - Area of single fin base, m$^2$}\\
\\
A_{fin} = 2wL_c
\text{\quad - Area of single fin, m$^2$}\\
\\
A_{fins} = NA_{fin}
\text{\quad - Area of all fins, m$^2$}\\
\\
A_{no fins} = 4L_{side}L_{duct}
\text{\quad - Area with fins removed, m$^2$}\\
\\
A_{unfin} = A_{nofins} - NA_b
\text{\quad - Area of fin bank with no fins, m$^2$}\\
\\
A_{tot} = A_{fins} + A_{unfin}
\text{\quad - Total area of fin bank, m$^2$}\\
\\
\\
\textbf{Fin Convection Resistance}\\
\\
R_{fins} = \frac{1}{hA_{fins}\eta_{fin}}
\text{\quad - Fins convection thermal resistance, K/W}\\
\\
R_{unfin} = \frac{1}{hA_{unfin}}
\text{\quad - Unfin convection thermal resistance, K/W}\\
\\
R_{tot} = \frac{1}{\frac{1}{R_{fins}} + \frac{1}{R_{unfin}}}
\text{\quad - Total resistance (parallel resistors), K/W}\\
\\
\textbf{(a) Determine Maximum allowable power dissipation per CPU}\\
\\
\dot{Q}_{tot} = \frac{T_b - T_\infty}{R_{tot}}
\text{\quad - Heat transfer rate, W}\\
\\
\dot{Q}_{tran} = \frac{\dot{Q}_{tot}}{N_{CPU}}
\text{\quad - Allowed power per CPU}\\
\\
\\
\textbf{(b) Determine required volume flow rate of air}\\
\\
\text{Assume perfectly caloric}\\
\\
\dot{m}_{air} = \frac{\dot{Q}_{tot}}{c_p\Delta T_{air}}
\)
\subsection{part2.m}
\(
P_{CPU} = 1000 \text{\quad - CPU Power, W}\\
P_{fan} = 1000 \text{\quad - Fan Power, W}\\
\\
P_{tot} = P_{CPU} + P_{fan}\\
\\
CPU_{perf} = 39000(1 - e^{-0.02P_{CPU}}) \text{\quad - CPU Performance Score}\\
\\
Q_{fan} = \frac{100000(\frac{P_{fan}}{10000})^{1/3}}{1000*60} \text{\quad - Fan Flowrate, m$^3$/s}\\
\\
\textbf{Thermal Resistances}\\
\\
R_{paste} = \frac{1}{h_cA_{c\ CPUs}}\\
\\
R_{cond} = \frac{t_{Duct}}{kA_{c\ wall}}\\
\\
\textbf{Mass Flow Rate}\\
\\
\dot{m}_{air} = Q_{fan}\rho\\
\\
\textbf{Mean Velocity}\\
\\
\text{Assume conservation of mass}\\
u_m = \frac{\dot{m}_{air}}{\rho A_c}\\
\\
\textbf{Surrogate Duct Reynolds Number}\\
\\
Re_D = \frac{u_m D_h}{\nu}\\
\\
\textbf{Convection Heat Transfer Coefficient}\\
\\
\text{Assume turbulent}\\
Nu = 0.023Re_D^{0.8}Pr^{0.4}\\
h = \frac{Nu*k}{D_h}
\text{\quad - Convection heat transfer coefficient (W/(m$^2$*k)}\\
\\
\textbf{Surface Temperatures}\\
\\
T_{\infty} = T_{bulk}\\
T_b = \frac{90R_{tot} + T_{\infty}(R_{paste} + R_{cond})}{R_{tot} + R_{paste} + R_{cond}}\\
\)
\subsection{constants.m}
\(
N = 24*4 \text{\quad - Total number of fins}\\
\\
N_{CPU} = 9*4 \text{\quad - Total number of CPUs}\\
\\
t = 0.002 \text{\quad - Fin thickness, meters}\\
\\
L = 0.025 \text{\quad Fin length, meters}\\
\\
L_{side} = 0.15 \text{\quad Duct side length, meters}\\
\\
L_{duct} = 0.15 \text{\quad Duct length, meters}\\
\\
w = L_{duct} \text{\quad Fin width, meters}
\)
\subsection{constants\_part1.m}
\(
\\
h = 50 \text{\quad Convection heat transfer coefficient, W/(m$^2$*k)}\\
\\
k = 178 \text{\quad - Thermal conductivity, W/(m*k)}\\
\\
T_\infty = 310 \text{\quad - Average air temp through box passageway, K}\\
\\
T_b = 360 \text{\quad - Maximum transistor temperature at base}\\
\\
c_p = 1006 \text{\quad - c$_p$ for air, J/(kg*K)}\\
\\
\rho = 1.14 \text{\quad - Air density, kg/m$^3$}\\
\\
\Delta T_{air} = 10 \text{\quad - Maximum change in air temp, K}
\)
\subsection{constants\_part2\_3.m}
\(
\\
t_{duct} = 0.01 \text{\quad - Duct thickness, meters}\\
\\
s = \frac{L_{side} - L - \frac{N}{4}t}{\frac{N}{4} + 1} \text{\quad - Spacing between fins, meters}\\
\\
A_c = L_{side}^2 - LtN \text{\quad - Cross sectional area, meters$^2$}\\
\\
D_h = \frac{2Ls}{L + s} \text{\quad - Hydraulic diameter of surrogate duct, meters}\\
\\
T_{mi} = 20 + 273.15 \text{\quad - Mean inlet temperature, K}\\
\\
T_{me} = 50 + 273.15 \text{\quad - Maximum mean duct exit temperature, K}\\
\\
c_p = 1007 \text{\quad - Specific heat of air between 15 and 70 deg C, J/(k*K) FROM TABLE}\\
\\
k = 383 \text{\quad - Thermal conductivity}\\
\\
A_{c\ CPUs} = 0.0014*9 \text{\quad - Total cross sectional area of CPUs for contact resistance, m$^2$}\\
\\
h_c = 6 \text{\quad - Thermal paste thermal contact coefficient, W/(m$^2$*K)}\\
\)
\end{document}
